\slajdid{09_2-s029}
\begin{frame}[label=09_2-s029]
\gusttekst\setlength{\abovedisplayskip}{3pt}\setlength{\belowdisplayskip}{3pt}\setlength{\abovedisplayshortskip}{2pt}\setlength{\belowdisplayshortskip}{2pt}Daljim porastom ulaznog napona, napon na drejnu tranzistora N opada, napon između sorsa i drejna tranzistora P raste, pa je moguće da i P tranzistor uđe u zasićenje.
\begin{columns}[c,onlytextwidth]\column{.21\textwidth}\centering\input{slike/tikz/s029_cmos_invertor.tex}\column{.77\textwidth}Uslov da tranzistor N radi u zasićenju
\[V_{DS,N} \ge \frac{(V_{GS,N}-V_{Tn})}{1+\frac{(V_{GS,N}-V_{Tn})}{L_nE_{Cn}}}\qquad V_O \ge \frac{(V_I-V_{Tn})}{1+\frac{(V_I-V_{Tn})}{L_nE_{Cn}}}\]
Uslov da tranzistor P radi u zasićenju
\[V_{DS,P} \le \frac{(V_{GS,P}-V_{Tp})}{1+\frac{(V_{GS,P}-V_{Tp})}{L_pE_{Cp}}}\qquad V_O-V_{DD} \le \frac{(V_I-V_{DD}-V_{Tp})}{1+\frac{(V_I-V_{DD}-V_{Tp})}{L_pE_{Cp}}}\]
\[V_O \le V_{DD}+\frac{(V_I-V_{DD}-V_{Tp})}{1+\frac{(V_I-V_{DD}-V_{Tp})}{L_pE_{Cp}}}\]\end{columns}
\end{frame}
